\section{Introduction}\label{sec:intro}

\subsection{Background}

A group of people at a wedding, a field trip or a hackathon takes hundreds of photos, and every photo ends up on the one phone that took it. Getting the photos to everyone else is still a manual job: someone creates a shared album days later, people forward pictures one by one in a chat, or the group gives up. The photos exist; the network between the phones is what is missing.

PHOTON closes that gap. Members join a live \emph{session} with a short code or a QR code, and from then on every photo any member takes appears in every other member's gallery, without anyone pressing a button. The server keeps a photo only for 30 minutes, long enough for the group's phones to fetch it.

Behind this simple behaviour sit most of the topics of a computer networks course:

\begin{itemize}
  \item \textbf{Application-layer protocols.} The system uses HTTP/1.1, HTTP/2 and HTTP/3 for requests and bulk transfer, WebSocket for a long-lived event channel, and push services (FCM, APNs) as wake-up signals.
  \item \textbf{Transport protocols.} TCP carries the control and event planes; QUIC over UDP carries photo data, because it survives a phone moving from Wi-Fi to cellular and it does not stall every stream when one packet is lost.
  \item \textbf{Reliable delivery over an unreliable network.} Large uploads are cut into checksummed chunks, sent several at a time, and only the missing chunks are resent, which is selective repeat at the application layer. Every request is idempotent, so a retry after a timeout is always safe.
  \item \textbf{Publish/subscribe with a durable log.} Events go into a Redis Stream and are read by one loop per client, so a client that reconnects receives exactly what it missed.
  \item \textbf{Network security.} TLS 1.3 protects every hop. Users sign in with OAuth 2.0 and PKCE and receive signed JSON Web Tokens. Photos are content-addressed by SHA-256, and the server recomputes the hash itself instead of trusting the client.
  \item \textbf{Distributed state.} Sessions and photo status live in Redis and change only through atomic server-side scripts. Late messages from a slow client are rejected with \emph{fencing tokens}, so they can never undo newer state.
\end{itemize}

\subsection{Motivation}

Four concrete problems drove the design of PHOTON.

\begin{enumerate}
  \item \textbf{Sharing is manual and late.} Shared albums and chat groups need someone to pick photos and send them. By the time that happens the event is over, and many photos are never shared.
  \item \textbf{Cloud albums keep everything, forever.} Existing services store every shared photo indefinitely, tie the group to one vendor's ecosystem, or recompress the images. A group that only wants to exchange photos during an event should not leave a permanent copy on someone else's server.
  \item \textbf{Mobile networks are hostile to large uploads.} A 4\,MB photo uploaded over a single HTTP request restarts from zero when the phone changes network, a lift cuts the signal or the operating system suspends the app. Retries that are not idempotent create duplicates, and notifications that are lost leave members with gaps.
  \item \textbf{A client's claim is not proof.} If the server trusts the hash a client announces, one member can upload different bytes under the hash of someone else's real photo, and deduplication then hides the real photo from everyone. The server has to verify the bytes before telling anyone the photo exists.
\end{enumerate}

\subsection{Application Domain}

\begin{pbox}{blue}{Where PHOTON Is Used}
\textbf{1. Events and trips:} weddings, college fests, family trips and sports days, where many people photograph the same moments and everyone wants everyone's photos.\par\smallskip
\textbf{2. Education and teams:} field labs, site visits and hackathons, where a team documents its work on several phones and needs one shared record during the activity.\par\smallskip
\textbf{3. Inspections and incident reporting:} surveyors, insurance assessors or volunteers who must get photos from the field to their colleagues at once, without keeping them on a third-party cloud afterwards.\par\smallskip
\textbf{4. Any short-lived group:} a session lasts hours, not years, and server copies are deleted 30 minutes after they are stored, so nothing outlives the event on the server.
\end{pbox}

\subsection{Problem Statement}

\begin{pbox}{red}{Core Problem Statement}
How can a group of phones, connected through unreliable and changing mobile networks, share every photo taken during a live session with every other member automatically, so that each photo is uploaded once, verified before anyone is told about it, delivered to every member without gaps or duplicates even across disconnections, and removed from the server within a bounded time?
\end{pbox}

\subsection{Objectives and Scope}

The objectives of the project were to:

\begin{enumerate}
  \item detect a new camera photo on the device and sync it while the user is in a session, also while the app is in the background;
  \item upload each photo once per session, using its SHA-256 hash as the content address;
  \item show a thumbnail first, then a preview, then save the original, according to network conditions;
  \item announce a photo to the group only after the server has verified its hash and stored it;
  \item deliver every event to every member at least once and in order, with deduplication on the client so that each photo is processed exactly once;
  \item make a photo unavailable exactly 30 minutes after it is stored and delete the file within 5 minutes after that;
  \item protect every hop with TLS 1.3 and every request with a verified identity.
\end{enumerate}

Out of scope for this version are video and Live Photos, photos taken before joining, long-term backup, more than one active session per user, and a web client (a browser cannot watch the phone's gallery).
